% ===== UNIT 31: VeRi-Wild: AQE, training-free fix (5.10 part 3) =====
% Status: REWRITTEN  (change to REWRITTEN / REVIEWED when done; see CHECKLIST.md)
% Rules: style rewrite only. Keep every \cite, \ref, \label, number, equation and table value as is. Change the wording, not the content.
Third, AQE helps the transferring stream at every scale ($+4.3$, $+3.7$ and $+2.9$\,pp mAP). It is an $\mathcal{O}(N)$ operation. It consistently hurts the Rank-1 of the weak stream, which is a known failure mode of query expansion on poorly clustered features. Fourth, $k$-reciprocal re-ranking reduces accuracy under transfer. In-domain it is the largest single accuracy step ($+5.55$\,pp on VeRi-776). Measured at the 3{,}000 protocol, it lowers ViT-B from 35.9 to 34.8\% mAP (and Rank-1 by 3.3\,pp), and the fusion from 33.6 to 30.9\%. On poorly clustered transfer features, the $k$-reciprocal neighbourhoods admit false matches, and the Jaccard blending then amplifies them. So the in-domain post-processing recipe does not carry over as a whole. Like the fusion weight, the benefit of re-ranking is domain-specific. This is not caused by the default parameters. We swept the same axis-aligned grid as in-domain (Section~\ref{sec:holdout}), and every re-ranked transfer configuration lands at 33.4--35.5\% mAP, below the 35.9\% of AQE alone. AQE itself is near-optimal at the VeRi-776-tuned $k{=}3$. The best value on the transfer side is $k{=}5$ (36.2\%), only $0.3$\,pp higher, and we deliberately do not adopt target-tuned settings. The re-ranking step costs 254\,s for the 41{,}861-image query$+$gallery set on CPU. This is roughly five times the 50\,s of the in-domain pass, for $3.2\times$ as many images (Section~\ref{sec:cost} discusses the scaling). We release per-query AP dumps for every transfer configuration alongside the VeRi-776 dumps (\texttt{scripts/eval/vw\_per\_query\_ap.npz}), so all reported CIs can be recomputed independently.

The fusion failure, unlike the re-ranking failure, has a training-free remedy. We tested query-adaptive weights computed from each stream's own AQE'd score row, with no target-domain information. One is a margin rule, $w_1^{(i)}{=}m_1/(m_1{+}m_2)$ with $m_s{=}\max(s)-\mathrm{mean}(\text{top-}50)$. The other is a top-100 dispersion analogue. At the 3{,}000 protocol the margin rule recovers \textbf{35.9\%} mAP (from 32.8\% at the fixed weight, and 35.9\% for the single stream) by shifting weight toward the reliable stream ($\bar{w}_1{\approx}0.63$). A constant-weight control explains the rescue. The transfer weight surface peaks at $w_1{\approx}0.8$/36.5\%, marginally above the single stream. So some complementarity survives the shift, but finding that optimum requires target-domain evaluation. A constant $w_1{=}0.63$ scores 36.0\%, and recentring the weights of the rule on the optimum adds nothing (36.4\%). So the value of the rule is finding a safe operating point without labels, not yet its per-query variation. In-domain the rules cost 1.1--1.4\,pp compared with the tuned fixed weight (87.0--87.3 vs.\ 88.4\% mAP at matched AQE-only settings), while staying well above the single stream (84.2\%). So adaptive weighting acts as a domain-robustness guard. A ground-truth per-query oracle reaches 40.1\% mAP / 66.9\% Rank-1, which is above every published direct-transfer entry in Table~\ref{tab:veriwild}. So the per-query headroom is real, but first-order score statistics do not capture it. Learned query-adaptive fusion is therefore a concrete, quantified direction for future work.
